\section{Comparison with Other Ensemble Methods}
\label{sec:comparison}

\subsection{The Five-Method Landscape}

The redistricting toolkit now contains five distinct search and sampling
methods.
Two are optimisation tools (they find good plans, not representative plans);
two are perturbation/exploration tools (they explore the neighbourhood of a
reference plan); and one --- \smc{} --- is a calibrated sampler (it draws
plans from a known distribution).
Table~\ref{tab:comparison} summarises the key dimensions.

\begin{table}[ht]
\centering
\caption{Comparison of redistricting ensemble and search methods.
  All empirical values are approximate and state-dependent.
  $n$ = census tracts, $k$ = districts, $T$ = convergence threshold
  (default 600 for \cs{}), $N$ = particle count (default 5{,}000 for \smc{}).
  ``Target distribution'' is what the method samples from, if anything.}
\label{tab:comparison}
\renewcommand{\arraystretch}{1.35}
\begin{tabular}{@{}p{2.5cm}p{2.8cm}p{1.6cm}p{2.4cm}p{3.5cm}@{}}
\toprule
Method & Target distribution & Mixing req.? & Per-plan cost & Primary use case \\
\midrule
\smc{}
  & Uniform over $\Pi$ (calibrated)
  & No
  & $O(k \cdot n \log n)$
  & Statistical inference ($p$-values) \\
\recom{}
  & Approx.\ uniform (Markov)
  & Yes
  & $O(T_{\rm mix} \cdot n/k \cdot \log n/k)$
  & Ensemble exploration \\
\textsc{ShortBurst}~\citep{cannon2022,g6paper}
  & None (optimisation)
  & N/A
  & $O(B \cdot L \cdot n/k \cdot \log n/k)$
  & Seat-outcome optimisation \\
\be{}~\citep{h1paper}
  & Conditional on structure
  & Within-node
  & $O({\rm nodes} \times n/k \cdot \log n/k)$
  & Hierarchical structure audit \\
\cs{}~\citep{b11paper}
  & None (deterministic)
  & N/A
  & $O(T \cdot n \log n / k)$
  & Certification / plan generation \\
\flip{}~\citep{g8paper}
  & None (local walk)
  & N/A
  & $O(F \cdot \deg_{\max})$
  & Local sensitivity analysis \\
\bottomrule
\end{tabular}
\end{table}

\subsection{SMC vs.\ ReCom: The Calibration Distinction}

Both \smc{} and \recom{} are designed to approximate the uniform distribution
over valid plans, but they differ fundamentally in how they achieve this.

\recom{} is a Markov chain: it defines a transition kernel $P(\pi, \pi')$ and
relies on $P$ having the uniform distribution as its stationary distribution.
Under this design, after a sufficient number of steps $T_{\rm mix}$, the
distribution of the current plan approximates the uniform distribution.
The difficulty is certifying that $T_{\rm mix}$ has been reached.
For redistricting chains on large state graphs, mixing is not guaranteed and
is rarely verified empirically before making inferential claims.

\smc{} is an importance sampler: it assigns weights $w_i$ to each generated
plan such that the weighted sample correctly represents the uniform distribution
by construction (Proposition~\ref{prop:weight-correct}).
No mixing assumption is required; the correction is analytic.

\begin{remark}[When ReCom claims are valid]
\recom{}-based $p$-value claims are statistically valid \emph{if} the chain
has mixed.
The practical challenge is that mixing time grows with state size and district
count; for California ($k = 52$) or Texas ($k = 38$), no published analysis
has demonstrated mixing for chain lengths feasible in practice.
\smc{} provides a valid alternative: for the same computational budget,
$N$ particles from \smc{} give a correct (though approximate) sample from the
uniform distribution, while $N$ steps of \recom{} from a fixed starting plan
give a walk of unknown distributional accuracy.
\end{remark}

\subsection{SMC vs.\ ConvergenceSweep: Inference vs.\ Certification}

\cs{}~\citep{b11paper} runs $T = 600$ independent METIS seeds and selects the
minimum-edge-cut plan.
The plans it generates are biased toward compact plans by design; the goal is
to find the best plan, not a representative sample.

\smc{} is designed for the opposite purpose: to generate a sample from which
inferential conclusions about the full plan space can be drawn.
The two methods are complementary, not competing:

\begin{itemize}
  \item Use \cs{} to generate the submitted plan.
  \item Use \smc{} to compute the $p$-value of the submitted plan relative to
        the uniform distribution.
\end{itemize}

The combination is powerful: \cs{} guarantees a good plan, and \smc{}
provides the court-facing statement that ``the submitted plan is at the
$X$th percentile of all valid plans by partisan outcome.''

\subsection{SMC vs.\ BisectionEnsemble: Global vs.\ Conditional}

\be{}~\citep{h1paper} samples plans that share the same hierarchical
bisection structure as the reference plan, varying only the sub-district
assignments within each node of the bisection tree.
This makes \be{} a \emph{conditional} ensemble: the distribution it samples
from is the uniform distribution over plans with a fixed hierarchical structure.

\smc{} samples from the \emph{unconditional} uniform distribution over all
valid plans, without fixing any hierarchical structure.
These serve different inferential purposes:
\begin{itemize}
  \item \be{} answers: ``Is there another way to draw districts with this
        general shape that produces a different partisan outcome?''
  \item \smc{} answers: ``What fraction of \emph{all} valid plans produce
        a given partisan outcome?''
\end{itemize}

For a full evidentiary picture, a practitioner may wish to report both: the
\smc{} global percentile establishes the plan's position in the full space,
while the \be{} conditional ensemble shows that the hierarchical structure
itself is not uniquely favourable.

\subsection{Computational Cost Comparison}

For North Carolina ($n \approx 2{,}200$, $k = 14$), approximate wall-clock
times on an 8-core workstation are:
\begin{itemize}
  \item \smc{} ($N = 5{,}000$ particles, Rayon): approximately 3--8 minutes.
  \item \recom{} ($T = 10{,}000$ steps): approximately 10--30 minutes
        (chain length is typically insufficient for mixing at $k = 14$).
  \item \cs{} ($T = 600$ seeds): approximately 3--8 minutes.
  \item \flip{} ($F = 10{,}000$ steps): approximately 1--3 seconds.
\end{itemize}

For inference, \smc{} is the correct choice regardless of cost: \recom{} at
any chain length does not provide a calibration guarantee, so spending more
time on \recom{} does not improve the validity of $p$-value claims.
\smc{} at $N = 5{,}000$ provides a valid (if approximate) inference with
an \ess{} of at least $\tau \cdot N = 2{,}500$ effective particles after
resampling.
